% =====================================================================
%  M1 - Section 2.1.1  Requirements Engineering for ML  (Team Vera)
%  Standalone preview for Overleaf. To merge into the Atenea template,
%  keep only the text between CONTENT START / CONTENT END and adjust the
%  heading levels to your template's numbering.
% =====================================================================
\documentclass[11pt]{article}
\usepackage[utf8]{inputenc}
\usepackage[T1]{fontenc}
\usepackage[margin=2.3cm]{geometry}
\usepackage{booktabs}
\usepackage{tabularx}
\usepackage{array}
\usepackage{enumitem}
\usepackage{amssymb}
\usepackage{hyperref}
\newcolumntype{Y}{>{\raggedright\arraybackslash}X}
\setlength{\tabcolsep}{5pt}
\renewcommand{\arraystretch}{1.18}

\begin{document}
% ==================== CONTENT START ====================
\section*{2.1.1\quad Selection of Problem and Requirements Engineering for ML}



\paragraph{Problem selection.}
Modern text-to-image generators (Stable Diffusion, DALL-E, MidJourney, \dots) produce photorealistic images at scale, fuelling misinformation, fake reviews, and trust problems for platforms and media. Reliably telling a \textbf{real photograph} from an \textbf{AI-generated image} cannot be solved with hand-written rules, the cues are subtle, generator-specific, and evolving, which makes it a genuine machine-learning problem. We target real-vs-AI detection in particular because nowadays it is increasingly difficult to recognise whether a piece of media is real or AI-generated.

\smallskip
\paragraph{ML goal.} Build, deploy, and operate an ML component that, given an input image, decides whether it is real or AI-generated, exposed as a monitored production service. Our main task is binary image classification(real vs ai\_generated), and optionally we might as well classify AI generated images among the 5 generators(which generator model owns the image).
Our focus is mostly on engineering around the model (reproducibility, QA, deployment, monitoring), not maximizing raw accuracy.

\medskip



\paragraph{What makes it unique.} Unlike general image-classification tasks whose target classes are usually obvious to a human (weather, animals, traffic signs, \dots), we target a category that in the near future may be very hard to classify even for people --- telling a real image from an AI-generated one. Our component also covers \textbf{five recent} generators (2022--2024), and instead of being specific to one dataset, our component would be easily configured and used with different, bigger datasets.




\paragraph{Intended use.} As main usecase of our ML component we can mention reliable platforms(such as news, social media, ...) to distinguish real from fake(or generated) images to prevent misinformation and correctly flag likely AI-generated uploads via an API, another usecase could be to quickly check the authenticity of an image with a confidence score, which is a well-documented and reproducible component for downstream developers to integrate in a larger system.


\subsection*{3.\quad ML component specification (I/O contract)}
\begin{itemize}[nosep,leftmargin=1.4em]
  \item \textbf{Input:} one image (JPEG/PNG/WEBP), bounded file size; preprocessed identically to training (fixed square resize + uniform re-encode).
  \item \textbf{Output:} JSON \texttt{\{"label": "real" | "ai\_generated", "score": <float 0--1>\}}; optionally generator model as well.
  \item \textbf{Operating context:} stateless REST service, containerised, with health and metrics endpoints.
\end{itemize}

\newpage

\subsection*{4.\quad Requirements}
\noindent\textit{Priority:} \textbf{M} = Must, \textbf{S} = Should, \textbf{C} = Could.

\subsubsection*{4.1\quad Functional (FR)}
\begin{tabularx}{\textwidth}{@{}l Y c@{}}
\toprule
\textbf{ID} & \textbf{Requirement} & \textbf{Pri} \\
\midrule
FR-1 & Classify a single image as \texttt{real} / \texttt{ai\_generated} & M \\
FR-2 & Return a confidence score in [0,1] with each prediction & M \\
FR-3 & Expose prediction via a REST endpoint (\texttt{POST /predict}) returning JSON & M \\
FR-4 & Accept JPEG/PNG/WEBP within a size limit; reject invalid input with a clear error & M \\
FR-5 & Apply the \textbf{same} preprocessing at train and serve time (resize + re-encode) & M \\
FR-6 & Provide health/readiness and metrics endpoints & S \\
FR-7 & Optionally return the most-likely generator (attribution) & C \\
\bottomrule
\end{tabularx}

\subsubsection*{4.2\quad Data (DR)}
\begin{tabularx}{\textwidth}{@{}l Y c@{}}
\toprule
\textbf{ID} & \textbf{Requirement} & \textbf{Pri} \\
\midrule
DR-1 & Use labelled real \& AI images from multiple recent generators & M \\
DR-2 & Data \& train/validation/test splits to be versioned and reproducible (DVC) & M \\
DR-3 & Handle class imbalance if observed in the data (via class weights or generator-balanced sampling) & M \\
DR-4 & Make sure no leakage such as image caption shared across train/val/test sets & M \\
DR-5 & Preprocessing to neutralise data to ensure the model does not rely on non-semantic \textbf{shortcuts} (such as aspect ratio, format) & M \\
DR-6 & Data usage complies with licence (CC BY 4.0, MS COCO) & M \\
DR-7 & Reserve a held-out generator (and data records) for drift evaluation & C \\
DR-8 & Automated data-quality checks (schema, class values, missingness, ...) & S \\
\bottomrule
\end{tabularx}

\subsubsection*{4.3\quad Model / performance (MR) --- \textit{targets (confirm, CI/CD?)}}
\begin{tabularx}{\textwidth}{@{}l Y c@{}}
\toprule
\textbf{ID} & \textbf{Requirement} & \textbf{Pri} \\
\midrule
MR-1 & Balanced accuracy (target $\geq$ 0.85) & S \\
MR-2 & Macro-F1 and recall on \textbf{minority} class, usually real images (targets $\geq$ 0.85 / $\geq$ 0.80) & S \\
MR-3 & \textbf{PR-AUC} (target $\geq$ 0.90) & S \\
MR-4 & \textbf{Cross-generator} (leave-one-generator-out) balanced accuracy (target $\geq$ 0.70) & S \\
MR-5 & Decision threshold based on validation results (not fixed 0.5) & M \\
MR-6 & Model retraining is reproducible (fixed seeds, versions, tracked data) & M \\
\bottomrule
\end{tabularx}

\smallskip
\noindent{Accuracy alone is misleading if there is class imbalance in data, so success would be defined on balanced/minority-aware metrics and on generalisation to unseen generators.}

\subsubsection*{4.4\quad Quality / non-functional (QR)}
\begin{tabularx}{\textwidth}{@{}l Y c@{}}
\toprule
\textbf{ID} & \textbf{Requirement} & \textbf{Pri} \\
\midrule
QR-1 & Image retrieval and processing runnable on general-purpose system resources (CPU, modest RAM) within a reasonable time (performance) & S \\
QR-2 & Stateless, containerised, and can handle multiple requests at the same time (scalability) & S \\
QR-3 & Code passes quality checks (Pylint/flake8), notebook lint (Pynblint), and unit tests (Pytest) (maintainability) & M \\
QR-4 & Track training energy use and CO$_2$ emissions using CodeCarbon (energy efficiency) & S \\
QR-5 & Monitor system resources and detect data/model drift (observability) & M \\
QR-6 & Health checks, proper handling of invalid input, best-effort uptime on free infra (reliability) & M \\
QR-7 & Validate inputs, and do not persist user images beyond the request (security/privacy) & M \\
QR-8 & Reproducible Docker image on the chosen platform (portability) & M \\
QR-9 & Automated pipeline that lints, tests, and builds/deploys the component on each change, e.g.\ GitHub Actions (CI/CD) & M \\
\bottomrule
\end{tabularx}

\subsubsection*{4.5\quad Challenges}
The project runs on general-purpose laptops and free-tier cloud, so we favour transfer learning on a small model, and the dataset must fit a free-tier DVC remote ($\sim$15\,GB) --- which Defactify ($\sim$7.5\,GB) does. We use the data under \textbf{CC BY 4.0} (non-commercial academic use,MS COCO). Assuming inputs resemble the training distribution (everyday photos vs full AI images) rather than face-swaps or adversarial inputs. The main risks are:


\medskip


\begin{itemize}[nosep,leftmargin=1.4em]
  \item \textbf{shortcut learning} --- the model keying on non-semantic cues such as image shape or format instead of ``AI-ness'' (mitigated by normalising size and re-encoding uniformly);
  \item \textbf{class imbalance} of roughly 5:1 fake:real (class weights / generator-stratified sampling and balanced metrics);
  \item \textbf{split leakage} from captions shared across train/val/test, which we already observed in EDA (a caption-grouped re-split);
  \item \textbf{generator drift}, i.e.\ degraded accuracy on unseen or newer generators (leave-one-generator-out evaluation and drift monitoring).
\end{itemize}

\paragraph{Success criteria} The component is successful when it is served via a \textbf{reproducible, containerised API}, and meets the model targets on a leakage-free test set with shortcuts controlled, while all processes monitored for drift.


\end{document}
